\section{Related work}\label{sec:related}

\subsection{Conformal prediction for time series}

Conformal prediction \citep{vovk2005algorithmic} turns any point predictor into a prediction set whose coverage is guaranteed in finite samples when the data are exchangeable. In the split variant, residuals on a held-out calibration set define a quantile, and the interval is the point forecast plus or minus that quantile. Time series are not exchangeable, and the guarantee does not hold in general. \citet{barber2023beyond} give a coverage bound beyond exchangeability, using weights that reflect how far the data depart from it. \citet{xu2021enbpi} propose a conformal interval for dynamic time series built on ensembles. \citet{zaffran2022agaci} study adaptive conformal predictions for time series and the role of the learning rate. A simple practical response is to compute the quantile only from the most recent residuals, which we call rolling conformal.

\subsection{Adaptive methods: ACI, FACI, conformal PID and SAOCP}

Adaptive conformal inference (ACI; \citealp{gibbs2021aci}) treats the miscoverage level as a quantity updated online. After each outcome, the working level moves by a step size times the difference between the target miscoverage and the realised miss indicator, so that the long-run frequency of misses approaches the target level without assumptions on the data. The step size is a tuning choice. FACI \citep{gibbs2024faci} aggregates several step sizes to avoid choosing one. Conformal PID control \citep{angelopoulos2023pid} treats the problem as control of a quantile and combines proportional, integral and forecasting components. SAOCP \citep{bhatnagar2023saocp} uses strongly adaptive online learning. The present benchmark includes ACI, FACI and conformal PID. SAOCP is implemented in the released \texttt{dracp} package and was run in a reproduction check (Section~\ref{sec:prespec}) but is not part of the panel benchmark.

These guarantees concern long-run miss frequency. They do not guarantee that every country, every horizon and every window is calibrated, and they are not stated for the situation in which the outcome of an $h$-step forecast is revealed only $h$ months later. How that delay affects the adaptive updates is an empirical question that this paper addresses with a diagnostic (Section~\ref{sec:feedback}).

\subsection{Regime-aware calibration: DRACP}

The following description is as described by the authors of DRACP (Dynamic Regime-Aware Conformal Calibration; \citealp{oancea2026dracp}, version 2, read from the HTML version). DRACP keeps the split-conformal radius but replaces the equal weights on past errors by a weighted quantile. The weight of a past error is a product of four factors: a temporal decay; a density-ratio correction for covariate shift; a Gaussian kernel whose effective sample size is floored; and a regime-similarity term from a Gaussian mixture model. An online controller adjusts the significance level from the realised coverage, aggregating adaptive-conformal experts over several learning rates. The authors evaluate the method on 48 macroeconomic and financial series (28 HICP inflation series of the euro area and EU member states, 13 US macroeconomic and energy series and 7 daily financial series) against split, rolling, ACI, FACI, SAOCP and conformal PID. The evaluation is mainly one-step-ahead, with an additional multi-horizon evaluation at $h=3,6,12$ with direct forecasts and delayed feedback. The authors state that the advantage is specific to one-step-ahead forecasting and depends on the point forecaster. The description was compared on 2026-11-23 with the HTML version 2 of the paper (author Bogdan Oancea; version 1 submitted 2026-08-17, version 2 submitted 2026-10-19; Sections 3 and 5 to 7 and Appendices B and C were read, the other appendices were not). It agrees with the paper on the weight factors, the 48 series, the baselines and the two limitations quoted, with four points to keep in mind. (i) In the parts read, the paper gives no numeric value for the calibration window, the target window, the decay, the number of regimes or the grid of learning rates; the configuration values used in this study (Section~\ref{sec:intervals}) come from the released package and the reproduction script and are not values stated by the paper. (ii) The paper calls its default controller a self-tuning aggregation of ACI experts over several learning rates and states that it is a heuristic of that paper that differs from the published FACI. (iii) In the paper's tables SAOCP has the best interval score and intervals about 20 percent narrower, while DRACP has the mean coverage closest to nominal; its mean coverage on the 28 HICP series is 0.878. (iv) In the paper's multi-horizon evaluation the mean coverage of DRACP falls from 0.890 at $h=1$ to 0.887, 0.873 and 0.842 at $h=3,6,12$, and the authors write that at $h\ge 3$ a purely online controller dominates.

Over 48 series and 20 seeds, the authors report mean coverage of 0.890 for DRACP, 0.884 for ACI, 0.883 for FACI, 0.868 for SAOCP, 0.841 for rolling, 0.825 for conformal PID and 0.781 for split conformal at nominal 0.90 (as recorded in the project protocol). This benchmark uses the released package through its public streaming interface. Our setting differs from theirs in three respects: the point model is refitted on an expanding window instead of being fitted once, the target is the cumulative $h$-month log growth instead of the annual rate, and the economies are commodity-dependent. The present study is a transfer study of a published method to a different panel; it draws no conclusion about the method beyond the limitation that its authors state.

\subsection{Inflation forecast uncertainty in emerging markets}

\paragraph{Search that was done.} On 2026-11-23 we ran a limited search: one session, one general web search engine and the record pages of arXiv, Crossref, NBER, the Dallas Fed and the Geneva Graduate Institute; English only; no citation chasing and no review protocol. The queries, in paraphrase, were: conformal prediction for inflation forecasting; conformal prediction for macroeconomic forecasting in emerging economies; prediction intervals for inflation forecasts and their calibration; empirical coverage of inflation fan charts and density forecasts; Dallas Fed inflation forecasting in emerging markets; foundation models for macroeconomic time series; zero-shot foundation models for inflation with prediction intervals; commodity price shocks and inflation in commodity-exporting economies. A work was kept only if its bibliographic record could be opened and it applied conformal or other interval methods to inflation or macroeconomic series, evaluated the calibration of inflation interval or density forecasts, evaluated pretrained models on macroeconomic series, or studied inflation responses to commodity prices. Completeness is not claimed, and for some works only the abstract was read.

\paragraph{What was found.} \citet{sengupta2025fewnet} forecast CPI inflation of Brazil, Russia, India and China at 12 and 24 months with a neural network and derive conformal intervals from a weighted conformal technique with a sliding window; in the text we read they report interval widths and no empirical coverage. \citet{brahem2026tunisia} forecasts inflation in Tunisia at 1, 3 and 6 months with machine-learning models and uses conformal prediction intervals (abstract read only). In neither did we find a comparison of conformal methods by coverage. The calibration of inflation intervals and densities in advanced economies is an older topic: \citet{diebold1997spf} report several deficiencies in the density forecasts of inflation of the US Survey of Professional Forecasters, and \citet{wallis2003fancharts} and \citet{clements2004boe} evaluate the Bank of England's density forecasts of inflation. The Dallas Fed working paper mentioned in the project notes is \citet{duncan2019emerging} (Working Paper 338, 2018), a comparison of inflation forecasting models for 14 emerging market economies at horizons of 1 to 12 quarters, concerned with point forecasts and not with intervals. We found no study of the calibration of inflation prediction intervals across emerging markets. Intervals for inflation are commonly derived from the forecast model's own error distribution, and we compare conformal calibration with the native ARIMA and ETS intervals of the engine in a secondary, post hoc analysis (Section~\ref{sec:native}).

\subsection{Pretrained time-series models}

Chronos \citep{ansari2024chronos} is a family of pretrained time-series forecasting models that can be applied to a new series without fitting, which is called zero-shot use. We use the smallest Chronos-Bolt model as a point forecaster and put conformal intervals on top of it (Section~\ref{sec:forecasters}). The model is described by its model card \citep{amazon2024chronosbolt}, which states no cutoff date for the training data. For macroeconomic series, \citet{jetwiriyanon2025zeroshot} back-test Chronos, TimeGPT and Moirai without fine-tuning on macroeconomic indicators and report reasonable uncertainty estimates and a degradation during rapid economic shocks, and \citet{michael2026tsfm} benchmark six zero-shot foundation models on energy, traffic and financial data and report a trade-off between point accuracy and probabilistic reliability (abstracts read only). We found no study of the calibration of foundation-model prediction intervals for inflation, and we make no claim about foundation models beyond their use here as a base forecaster under conformal intervals.

\subsection{Positioning}

The paper is a transfer and benchmark study, not a new method. Its aim is to measure, with decision rules recorded in advance, whether the published ranking of conformal methods holds on a different panel, and to report null outcomes as such. The closest work on method is the DRACP paper \citep{oancea2026dracp}, and the closest applications of conformal intervals to inflation in emerging economies are \citet{sengupta2025fewnet} and \citet{brahem2026tunisia}. We do not replicate the setup of DRACP; we reproduce it only partially (Sections~\ref{sec:prespec} and~\ref{sec:limitations}). The novelty claim is limited to a benchmark of several adaptive conformal methods by empirical coverage, at several horizons with delayed outcomes, with decision rules recorded in advance; conformal intervals for inflation in emerging economies as such are not new.
